%% SCRAP: archive/operations/qa-intake-workflow
%% SOURCE: docs/working/archive/operations/qa-intake-workflow.adoc
%% STATUS: OBSOLETE
%% FITS: none
%% EDITORIAL: lifted — prose rewritten to press voice

\section{QA Intake and CAPA Backlog Workflow (Archived)}

This document described the quality assurance intake process for StarForth
at the time GitHub-based governance workflows were active (November 2025).
The workflow has since been superseded by the current governance model.
The record is preserved for process archaeology.

\subsection{Overview}

The QA intake workflow routed defects from five source types into the CAPA
(Corrective and Preventive Action) backlog:
manual testing, CI/CD failures, user reports, developer self-reports, and
code review findings. GitHub Actions validated submissions, applied severity
labels, and created Kanban board entries automatically.

\subsection{Defect Sources}

\begin{center}
\begin{tabular}{lll}
\toprule
Source & Description & Trigger \\
\midrule
Manual testing    & QA team testing          & QA creates issue via template \\
CI/CD failures    & Baseline torture test    & Jenkins auto-creates CAPA \\
User reports      & GitHub Issues            & QA converts to CAPA format \\
Developer reports & Developer self-report    & Developer creates directly \\
Code review       & Reviewer finding in PR   & Reviewer creates linked issue \\
\bottomrule
\end{tabular}
\end{center}

\subsection{CAPA Issue Structure}

Each CAPA issue required eight sections: problem description, reproduction
steps, expected behavior, code location (if known), severity assessment
(critical / major / minor / low), regression indicator, workaround, and
a minimal test case in FORTH.

\subsection{Kanban Flow}

Issues transitioned through five columns: Backlog, QA Review, In Progress,
In Review, and Done. CI/CD failures auto-created issues and placed them in
QA Review. Developer PRs referencing \texttt{Closes \#CAPA-NNN} moved cards
to In Progress. Merge to master auto-closed issues.

\subsection{Automated Severity Classification}

GitHub Actions auto-classified severity and applied labels
(\texttt{type:capa}, \texttt{severity:\{level\}}, optionally
\texttt{regression}) upon issue creation. FMEA requirement was determined
during QA triage: required for core arithmetic or safety-critical changes;
optional for localized bugs.

%% TODO(bob): confirm whether any successor governance process is worth
%% documenting in the dev-guide volume, or whether this is purely archival.
